%
% File: chapter03.tex
% Chapter 3: Results and Discussion
%
\chapter{Results and Discussion}
\label{chap:results}
\setlength{\parskip}{1em}

\section{Overall Performance Summary}

Table \ref{tab:results} presents the comprehensive performance comparison of all ten regression algorithms evaluated in this study. The results reveal substantial performance variation across different algorithmic families, with ensemble tree-based methods demonstrating clear superiority over linear regression techniques.

\begin{table}[H]
\centering
\caption{Comparative Performance of Regression Algorithms on Kazakhstan Divorce Statistics}
\label{tab:results}
\begin{tabular}{|l|c|c|c|c|c|}
\hline
\textbf{Algorithm} & \textbf{Features} & \textbf{Targets} & \textbf{k-fold} & \textbf{CV RMSE} & \textbf{CV R²} \\
\hline
Ridge Regression & 5 & 1 & 10 & 3267.06 $\pm$ 740.64 & -0.1810 $\pm$ 0.1763 \\
Lasso Regression & 5 & 1 & 10 & 3042.07 $\pm$ 815.24 & 0.0125 $\pm$ 0.0137 \\
Elastic Net & 5 & 1 & 10 & 3042.31 $\pm$ 822.44 & 0.0145 $\pm$ 0.0056 \\
KNN Regression & 5 & 1 & 10 & 3267.06 $\pm$ 740.64 & -0.1810 $\pm$ 0.1763 \\
\textbf{Extra Trees} & 5 & 1 & 10 & \textbf{376.16 $\pm$ 112.81} & \textbf{0.9831 $\pm$ 0.0083} \\
AdaBoost & 5 & 1 & 10 & 1100.18 $\pm$ 91.03 & 0.7915 $\pm$ 0.2711 \\
Gradient Boosting & 5 & 1 & 10 & 831.61 $\pm$ 159.13 & 0.9158 $\pm$ 0.0338 \\
XGBoost & 5 & 1 & 10 & 835.80 $\pm$ 165.52 & 0.9146 $\pm$ 0.0366 \\
\textbf{CatBoost} & 5 & 1 & 10 & \textbf{572.12 $\pm$ 73.27} & \textbf{0.9547 $\pm$ 0.0321} \\
HistGradientBoosting & 5 & 1 & 10 & 914.38 $\pm$ 177.86 & 0.8994 $\pm$ 0.0377 \\
\hline
\end{tabular}
\end{table}

\section{Performance Analysis by Algorithm Family}

\subsection{Linear Regression Models}

The linear regression family (Ridge, Lasso, Elastic Net) demonstrated poor performance on this dataset, with RMSE values exceeding 3000 and near-zero or negative R² scores. Ridge Regression achieved RMSE of 3267.06 and R² of -0.1810, indicating that the model performs worse than simply predicting the mean divorce count. The negative R² values signify that linear assumptions fundamentally fail to capture the underlying relationships in this data.

Lasso and Elastic Net showed marginal improvement with slightly positive R² values (0.0125 and 0.0145 respectively), suggesting minimal variance explanation despite L1 regularization's feature selection properties. The high RMSE values (approximately 3042) render these models impractical for real-world divorce prediction applications.

These results strongly indicate that divorce patterns across Kazakhstan's regions exhibit significant non-linear relationships that cannot be adequately modeled through linear combinations of the engineered features. The failure of regularization techniques suggests that the issue lies not in overfitting or multicollinearity, but rather in the fundamental model structure.

\subsection{K-Nearest Neighbors}

KNN Regression, despite feature scaling, performed identically to Ridge Regression (RMSE: 3267.06, R²: -0.1810). This unexpected equivalence likely results from the curse of dimensionality in the five-dimensional feature space combined with relatively sparse regional coverage. The default k=5 configuration may have proven suboptimal for the dataset's structure, where nearest neighbors in feature space do not necessarily correspond to similar divorce patterns.

The poor KNN performance also suggests that simple distance-based similarity in the engineered feature space does not effectively capture divorce trend similarities across regions. Geographical proximity and temporal adjacency may require more sophisticated representation than Euclidean distance in standardized feature space.

\subsection{Ensemble Tree-Based Methods: Superior Performance}

The most striking finding of this study is the exceptional performance of ensemble tree-based methods, with Extra Trees Regression achieving the best overall results.

\subsubsection{Extra Trees Regression: Best Overall Performance}

Extra Trees demonstrated outstanding predictive accuracy with RMSE of 376.16 $\pm$ 112.81 and R² of 0.9831 $\pm$ 0.0083. These results indicate that the model explains 98.31\% of variance in divorce counts, representing near-perfect prediction capability. The relatively low standard deviation (112.81) further indicates stable performance across different data subsets.

The success of Extra Trees can be attributed to several factors. First, the algorithm's extreme randomization in both feature selection and split point determination creates highly diverse base learners, reducing overfitting while maintaining strong individual tree performance. Second, the ensemble averaging process effectively smooths out noise while preserving the underlying non-linear patterns. Third, the algorithm's insensitivity to feature scaling eliminates potential preprocessing artifacts.

Notably, Extra Trees' training time (0.1535 seconds) remains reasonable despite the large ensemble size (100 trees), making it practical for operational deployment. The algorithm's ability to provide feature importance scores could offer valuable insights into which factors (district, year, area type) most strongly influence divorce rates, though such analysis falls outside the current study's scope.

\subsubsection{CatBoost: Second-Best Performance}

CatBoost achieved the second-best performance with RMSE of 572.12 $\pm$ 73.27 and R² of 0.9547 $\pm$ 0.0321, explaining 95.47\% of variance. While slightly less accurate than Extra Trees, CatBoost offers several advantages. Its specialized handling of categorical features through ordered target statistics may have better leveraged the encoded district information. The lower standard deviation (73.27) compared to other boosting methods suggests superior generalization stability.

CatBoost's competitive performance validates the effectiveness of modern boosting implementations specifically designed for structured data with categorical variables. The algorithm's built-in overfitting protection through ordered boosting and symmetric tree structures likely contributed to robust cross-validation performance.

\subsubsection{Gradient Boosting Methods: Strong Performance}

Gradient Boosting (RMSE: 831.61, R²: 0.9158) and XGBoost (RMSE: 835.80, R²: 0.9146) delivered nearly identical performance, both achieving R² values above 0.91. The similarity in results despite different implementations suggests that both algorithms have converged to similar underlying function approximations.

HistGradientBoosting showed competitive performance (RMSE: 914.38, R²: 0.8994), though slightly behind the other gradient boosting variants. The histogram-based approach trades some accuracy for computational efficiency, which may be valuable in production settings requiring frequent retraining.

\subsubsection{AdaBoost: Moderate Performance}

AdaBoost achieved moderate results (RMSE: 1100.18, R²: 0.7915) with notably high cross-validation variance ($\sigma_{R^2}$ = 0.2711). This variability suggests that AdaBoost's sequential weighting scheme may be sensitive to the particular sample distribution in each fold. While still outperforming linear methods by substantial margins, AdaBoost trails modern gradient boosting implementations, likely due to its reliance on simple weighted averaging rather than gradient-based optimization.

\section{Test Set Performance}

Test set evaluation confirmed the cross-validation findings, with Extra Trees achieving test RMSE of 325.39 and test R² of 0.9822, slightly better than cross-validation performance. This indicates excellent generalization without overfitting. CatBoost similarly maintained strong test performance (RMSE: 526.38, R²: 0.9534).

Conversely, linear models continued to underperform on the test set, with Ridge achieving test R² of -0.1201, confirming the systematic inadequacy of linear assumptions for this prediction task.

The consistency between cross-validation and test performance across all algorithms validates the robustness of the 10-fold cross-validation methodology and supports the reliability of the findings.

\section{Computational Efficiency}

Training time analysis reveals interesting trade-offs between accuracy and computational cost. Linear models train in under 0.01 seconds, offering near-instantaneous results despite poor accuracy. Extra Trees requires 0.1535 seconds, providing the best accuracy-to-time ratio. Among boosting methods, CatBoost (0.2125s) and XGBoost (0.2545s) require the most time, reflecting their sophisticated optimization procedures.

For operational deployment, Extra Trees offers an attractive balance: state-of-the-art accuracy with training times under 200 milliseconds, enabling frequent model updates as new divorce statistics become available.

\section{Implications for Divorce Prediction in Kazakhstan}

The findings carry several important implications for understanding and predicting divorce patterns in Kazakhstan:

\textbf{Non-linear Regional Dynamics:} The failure of linear models and success of ensemble methods strongly suggests that divorce patterns are influenced by complex, non-linear interactions between regional characteristics, area types, and temporal factors. Simple linear relationships between years, districts, and divorce counts do not exist; rather, the relationships are moderated by context-dependent factors.

\textbf{Regional Heterogeneity:} The superior performance of tree-based models indicates substantial heterogeneity across Kazakhstan's regions. Different districts likely respond differently to temporal trends, with urban and rural areas exhibiting distinct divorce dynamics. Tree-based models' ability to learn region-specific patterns explains their predictive advantage.

\textbf{Practical Application:} The high accuracy of Extra Trees and CatBoost enables practical applications in social planning. Government agencies could deploy these models to forecast regional divorce trends, informing resource allocation for family support services, legal aid programs, and counseling facilities. Predictions could be generated at district level with high confidence, allowing targeted interventions.

\textbf{Feature Engineering Opportunities:} The strong performance despite limited features (only 5) suggests that the basic temporal and geographical information captures fundamental drivers of divorce patterns. However, the remaining unexplained variance (approximately 1.69\% for Extra Trees) indicates potential for improvement through incorporation of socioeconomic variables such as employment rates, income levels, and educational attainment.

\section{Comparison with Existing Literature}

These findings align with recent literature demonstrating tree-based ensemble superiority for structured data prediction tasks. The results echo Fernández-Delgado et al.'s findings that random forest variants consistently rank among top performers, while extending this conclusion to divorce prediction in a post-Soviet context.

The magnitude of performance difference between linear and tree-based methods (R² difference exceeding 0.96) exceeds typical gaps observed in many machine learning benchmark studies, suggesting particularly strong non-linearity in the divorce statistics domain. This observation supports Borisov et al.'s assertion that tree-based models remain the gold standard for tabular data, especially when non-linear relationships dominate.

\section{Limitations and Considerations}

Several limitations warrant acknowledgment:

\textbf{Interpretability Trade-off:} While Extra Trees and CatBoost achieve superior accuracy, they operate as "black boxes" compared to linear models. Understanding why specific predictions are made requires additional interpretation techniques such as SHAP values or partial dependence plots.

\textbf{Temporal Generalization:} The models' ability to extrapolate beyond the 2000-2023 training period remains uncertain. Significant societal changes could render historical patterns less predictive of future trends.

\textbf{Aggregation Level:} The dataset represents regional aggregates rather than individual-level data. Micro-level patterns influencing individual divorce decisions may differ from macro-level regional trends captured by these models.

\textbf{Causality:} These models identify predictive patterns but do not establish causal mechanisms. The high predictive accuracy does not imply that year or district directly cause divorce; rather, they serve as proxies for underlying social, economic, and cultural factors.

\section{Synthesis}

This comprehensive comparative analysis definitively demonstrates that ensemble tree-based methods, particularly Extra Trees Regression and CatBoost, vastly outperform traditional linear regression approaches for Kazakhstan divorce statistics prediction. The results underscore the importance of algorithm selection based on data characteristics, with non-linear ensemble methods proving essential for social science applications exhibiting complex regional and temporal dynamics.

The findings provide both methodological guidance for researchers analyzing demographic trends and practical tools for policymakers seeking data-driven insights into divorce patterns across Kazakhstan's diverse regions.

\chapter{Conclusion}
\label{chap:conclusion}
\setlength{\parskip}{1em}

This study conducted a rigorous comparative analysis of ten machine learning regression algorithms applied to divorce statistics prediction in Kazakhstan, employing 10-fold cross-validation on 8,458 records spanning 2000-2023 across multiple regions. The research addressed a critical gap in understanding which algorithmic approaches yield optimal performance for demographic prediction in post-Soviet Central Asian contexts.

\section{Key Findings}

The investigation revealed several definitive conclusions:

\begin{enumerate}
    \item \textbf{Ensemble Superiority}: Tree-based ensemble methods dramatically outperformed linear regression techniques, with Extra Trees Regression achieving the highest accuracy (RMSE: 376.16, R²: 0.9831), followed by CatBoost (RMSE: 572.12, R²: 0.9547).
    
    \item \textbf{Linear Model Inadequacy}: Ridge, Lasso, Elastic Net, and KNN regression demonstrated poor performance with near-zero or negative R² values, indicating fundamental inability to capture non-linear divorce patterns.
    
    \item \textbf{Gradient Boosting Effectiveness}: Modern gradient boosting implementations (XGBoost, Gradient Boosting, HistGradientBoosting) achieved R² values exceeding 0.89, validating their utility for structured demographic data.
    
    \item \textbf{Non-linear Relationships}: The stark performance gap between linear and tree-based methods confirms complex, non-linear interactions between regional characteristics, urbanization, and temporal factors in determining divorce rates.
    
    \item \textbf{Practical Applicability}: The best-performing models achieve sufficient accuracy (R² > 0.95) for operational deployment in social policy planning and resource allocation.
\end{enumerate}

\section{Contributions}

This research makes several valuable contributions:

\textbf{Methodological}: Establishes best practices for machine learning application to divorce prediction in Central Asian contexts, providing a rigorous comparative framework applicable to similar demographic prediction tasks.

\textbf{Empirical}: Generates novel insights into divorce pattern complexity across Kazakhstan's regions, demonstrating the necessity of non-linear modeling approaches for this domain.

\textbf{Practical}: Delivers validated predictive models suitable for real-world application by government agencies and social service organizations in Kazakhstan.

\textbf{Theoretical}: Contributes to understanding of algorithm selection principles for social science prediction tasks, particularly in contexts characterized by regional heterogeneity and temporal dynamics.

\section{Recommendations for Practice}

Based on the findings, we recommend:

\begin{itemize}
    \item \textbf{Algorithm Selection}: Practitioners working with similar demographic datasets should prioritize ensemble tree-based methods (Extra Trees, CatBoost) over traditional linear regression approaches.
    
    \item \textbf{Model Deployment}: Kazakhstan government agencies should consider deploying Extra Trees or CatBoost models for operational divorce trend forecasting, enabling data-driven resource planning.
    
    \item \textbf{Feature Enhancement}: Future work should incorporate socioeconomic variables (employment, education, income) to potentially improve the already-high prediction accuracy.
    
    \item \textbf{Regular Retraining}: Given rapid societal changes, predictive models should be retrained annually as new divorce statistics become available.
\end{itemize}

\section{Future Research Directions}

Several promising avenues for future investigation emerge from this work:

\begin{enumerate}
    \item \textbf{Hyperparameter Optimization}: Systematic hyperparameter tuning through grid search or Bayesian optimization could potentially improve model performance beyond the baseline configurations employed in this study.
    
    \item \textbf{Feature Importance Analysis}: Detailed examination of feature contributions using SHAP values or permutation importance could reveal which factors most strongly drive divorce predictions.
    
    \item \textbf{Temporal Modeling}: Incorporating time-series specific methods (ARIMA, Prophet, LSTM networks) could better capture temporal dependencies and seasonal patterns.
    
    \item \textbf{Multi-regional Comparison}: Extending this methodology to other Central Asian nations (Uzbekistan, Kyrgyzstan, Tajikistan) would test generalizability and reveal region-specific patterns.
    
    \item \textbf{Causal Analysis}: Applying causal inference techniques to identify mechanistic drivers of divorce trends, moving beyond pure prediction to explanation.
    
    \item \textbf{Individual-level Modeling}: If individual-level data becomes available, comparing aggregate versus individual prediction performance could yield valuable methodological insights.
\end{enumerate}

\section{Closing Remarks}

This research demonstrates that careful algorithm selection and evaluation can yield highly accurate predictive models for complex social phenomena such as divorce patterns. The findings underscore the value of modern machine learning techniques for demographic research and social policy applications, particularly in under-researched Central Asian contexts. By systematically comparing ten diverse algorithms under rigorous evaluation protocols, this study provides both methodological guidance and practical tools for researchers and policymakers seeking to understand and predict family dynamics in Kazakhstan and similar societies undergoing rapid socioeconomic transformation.

The exceptional performance of ensemble tree-based methods reveals that divorce patterns are shaped by intricate, non-linear interactions between geographical, temporal, and social factors—patterns that simple linear models cannot capture but sophisticated machine learning algorithms can effectively learn. As Kazakhstan continues to evolve socially and economically, data-driven approaches like those demonstrated here will prove increasingly valuable for evidence-based policy formulation and social service planning.
